% !TeX root = ../main.tex
% !TeX program = xelatex
% !TeX encoding = UTF-8
\section{研究问题与技术贡献}

供热规划需要确定设备配置、逐时运行及费用与排放的取舍。余热要满足温度和输送条件，热泵要在低温时保持出力，多楼预热还会竞争区域用电容量，三者需要联合计算。

建筑热动态模型能够描述供热对后续室温的影响\supcite{2}；典型时段方法缩减计算规模，但需处理蓄热状态的时序联系\supcite{4}；$\varepsilon$约束方法通过改变目标上限比较方案\supcite{5}。本文重点研究这些方法在供热问题中的衔接：压缩后保留什么负荷特征，预热占用多少容量，以及方案能否满足设备和管网要求。

按赛事要求\supcite{1}，以乌鲁木齐、固原和武汉为背景，将热源、管网、建筑负荷与储热设备纳入源—网—荷—储模型，形成四项工作。

第一，统一热量计量。直接换热与热泵共用的余热只计一次，热泵源侧取热和电功分开计算；逐楼室温、储热量与系统电力共同约束运行。

第二，校正代表日权重。代表日是选出的实际日期，权重表示折合天数。在日数和能量守恒下调整权重，使负荷分布接近原序列，同时检查峰值和相邻小时的最大变化。

第三，将工程条件反馈给优化。厂家同工况制热量和输入功率同时限制出力与耗电；不满足远端温度或压降要求的管网模式被排除后重算，保存调整原因与费用。

第四，通过对照解释决策收益。完整服务年回放检验连续运行；同设备策略、热容与配电上限、碳帽和采购预算对照分别识别节费机制及方案切换条件。

\section{数据基础与三城实验设计}
\subsection{数据来源和用途}

模型输入包括天气、建筑参数、负荷记录、设备性能和费用因子。每份资料保留来源、所属对象与时间范围，计算得到的参数同时保留换算过程。表\ref{tab:data-scope}列出它们在实验中的用途，区分实测负荷、公共建筑原型和模型计算值。

\begin{table}[htbp]
\centering
\caption{研究数据及其用途}\label{tab:data-scope}
\begin{tabularx}{\textwidth}{@{}>{\raggedright\arraybackslash}p{0.23\textwidth}X >{\raggedright\arraybackslash}p{0.23\textwidth}@{}}
\toprule
数据来源 & 在本文中的作用 & 数据类型\\
\midrule
NASA POWER逐时气象\supcite{9} & 三城同一时间段的天气输入；统一时区与单位 & 公开气象网格数据\\
BESTEST 600原型\supcite{11} & 建筑几何与热工参数的来源；用于推导简化热模型参数 & 公共原型的派生参数\\
BDG2建筑计量\supcite{10} & 六栋建筑的负荷压缩与代表日权重对照 & 外部计量基准\\
TICA产品样本\supcite{12} & 两型号共165个制热量/功率点及缺测角点 & 厂家公开工况表\\
电价及排放因子\supcite{13,14} & 在共同计费条件下比较电量费与碳排放 & 公开价格和核算因子\\
OSM空间数据 & 三城建筑与道路候选的空间参考 & 公开空间数据\\
\bottomrule
\end{tabularx}
\end{table}

三城气候对照从北京时间2024年11月16日00:00开始，到2025年3月16日00:00结束，不含终点时刻，共120日、2880个连续小时。各城采用相同建筑和室温要求，仅更换天气输入。费用统一按湖北工商业分时电价计算，碳排放统一由用电量乘固定平均排放因子得到，以便比较天气和调度的影响。这里得到的是共同计费条件下的结果。厂家热泵实验由48~m$^2$公共原型按20倍规模构造960~m$^2$计算对象，允许室温在20--27~$^\circ$C内变化，开始与结束时均为20~$^\circ$C。

三城公共原型用于气候与调度对照，六栋外部计量建筑用于检验负荷压缩，分别报告结果。

\subsection{三类场景与验证分工}

乌鲁木齐关注工业余热和集中输配，固原关注低温设备与清洁替代，武汉关注冷热需求及多楼协同。三城参与同窗调度对照；网络反馈用72小时参考系统，热容交互用武汉双楼算例。另保留武汉四楼、两区域的历史数学场景：2023年7月1日至2024年7月1日的8784小时服务年，包含2023年11月15日至2024年3月16日的2928小时声明冬季。日期均为北京时间、含起点不含终点；该冬季是项目服务假定。

BDG2数据来自境外建筑，其实测负荷用于评价压缩误差。三城OSM数据提供建筑位置和候选道路路径；热源能力、供热对象与逐时出力则由模型配置确定，不能仅从道路或建筑轮廓推得。

表\ref{tab:experiment-design}列出各实验固定了什么、改变了什么，以及由此能够判断什么。

\begin{table}[htbp]
\centering
\caption{主要实验的控制变量与验证任务}\label{tab:experiment-design}
\begin{tabularx}{\textwidth}{@{}>{\raggedright\arraybackslash}p{0.19\textwidth}XX@{}}
\toprule
实验对象 & 保持一致的条件 & 改变因素及验证任务\\
\midrule
六栋计量建筑 & 观测序列、各楼实际日集合 & 权重方法；检验分布保真与统计准入\\
三城两型号设备 & 各组天气、整机配置与温度服务 & 调度策略；识别同资产运行收益\\
武汉四楼两区域 & 合成参数、既定资产及舒适要求 & 完整服务年及声明冬季回放；检验连续状态\\
武汉双建筑 & 传热系数、天气与设备 & 热容倍率及公用电上限；识别交互作用\\
72小时参考管网 & 交付需求与物理校核阈值 & 可用模式目录；量化反馈修正及费用\\
三城理想资产 & 各城天气、资产与服务条件 & 碳排上限；比较降费、碳排与峰值\\
条件投资案例 & 2880小时服务与性能假设 & 设备台数、采购预算；确定选择阈值\\
\bottomrule
\end{tabularx}
\end{table}

\subsection{统一比较与计量规则}

每组对照仅改变待研究的因素，其余条件按表\ref{tab:experiment-design}保持一致。功率、能量、费用与排放分别用kW、kWh、元与kg计量。代表日权重用于汇总费用和能量；检验室温或储热量能否连续变化时，仍按真实时间顺序计算。程序检查年份、时区及小时是否连续，并保存配置、文件校验码、求解状态和独立复算结果，以确定每个数值对应的输入与计算过程。
